\anachapter{De l'atome de lithium aux matériaux du pôle $\oplus$}
%\anachapter{Propriétés élèmentaires du lithium}
\label{chp:atom_to_mat}

\objectif{
\faBattery[0] De l'atome au matériau
\begin{itemize}[label=\faAngleRight]
\item Connaître la classification de dimensionnement et citer quelques matériaux d'intérêt dans les batteries ;
\item Savoir donner un sens thermodynamique sur la base d'arguments enthalpique et entropique justifiant l'insertion du lithium dans une structure ;
\item Avoir quelques idées sur les différents modes d'insertion du lithium dans une structure ;
\item Savoir interpréter la stabilité d'une phase/structure/variété allotropique à l'aide d'une courbe de solvus ;
\item Justifier l'évolution de la f.e.m. à l'aide de la stabilité des phases, du nombre de phase dans une courbe de charge/décharge.
\end{itemize}

\faBattery[2] Stabilité des structures d'insertion les plus courantes
\begin{itemize}[label=\faAngleRight]
\item Connaître les paramètres stabilisant une charpente moléculaire ;
\item Identifier les sites d'insertion dans une structure ;
\item Connaître les paramètres pouvant jouer un rôle dans la stabilité des sites d'insertion lors de l'occupation par des ions lithium.
\end{itemize}

\faBattery[4] Stabilité des liaisons iono-covalentes
\begin{itemize}[label=\faAngleRight]
\item Identifier la coordination des cations dans la charpente moléculaire et justifier de la stabilité de l'ion ;
\item Savoir juger de la stabilité d'une structure hôte en utilisant le modèle de liaison ionique, sur la base d'arguments structuraux (rayons, polyèdres, connectivité$\ldots$) et énergétiques (champs cristallin) ;
\item Interpréter des modifications structurales (\emph{par ex.} paramètres de maille) ou des paramètres énergétique (\emph{par ex.} f.e.m.) sur la base d'une interaction à deux orbitales.
\end{itemize}
}

%\section{De l'atome aux matériaux}

Dans cette section, nous allons aborder l'aspect structural des matériaux d'insertion du pôle $\oplus$. Comme souvent en chimie, les chimistes des matériaux disposent d'un pool restreint de "structures cristallographiques" qui peuvent être viables pour les batteries. \`{A} partir de ce pool, les chimistes réalisent des substitutions d'éléments afin d'améliorer les performances d'un matériau et de répondre au cahier des charges fixé.

%WARNING ci-dessous
%\ce{M^{no} O2 + xLi^+ + x e^- = A_x M^{no-x} O2}
%Dualité ion-électron = ion associé à la structure cristallographique

Dans la majorité des technologies de batteries d'aujourd'hui, il est important de se souvenir que les phénomènes ayant lieu \blancbis{au sein} du matériau, sont tout aussi importants que les phénomènes ayant lieu à l'interface du matériau et de l'électrolyte\mySideNote{L'étude de la réactivité, des phenomènes d'adsorption$\ldots$ ayant lieu à l'interface des solides est un des principaux axes de recherche en électrochimie et sera abordé dans le module associé.}.

De ce fait, il est important de penser la structure des matériaux comme quelque chose de \blancbis{dynamique}, et non pas de statique et immuable. En effet, lorsqu'il y a insertion du lithium dans un matériau, plusieurs questions se posent : 
\begin{itemize}
\item le lithium stabilise-t-il ou destabilise-t-il le matériau ? 
\item dans quels sites cristallographiques le lithium s'insère-t-il de manière préférentielle ?
\item le matériau hôte est-il stable ? Si oui, va-t-il ne connaître qu'une légère dilatation de sa structure ? Sinon, va-t-il subir un changement de structure (transition de phase) ou carrément se vitrifier ?
\end{itemize}

\section{Le lithium au sein d'un matériau hôte}

\subsection{Topologie des matériaux d'insertion : la dimension des matériaux}

Les matériaux d'insertion sont des structures disposants de sites interstitiels\mySideNote{En cristallographie, vous avez vu jusque-là les sites octaédriques et les sites tétraédriques.} pouvant accueillir des ions lithium.

\begin{figure}[H]
\begin{centering}
\subfloat[\label{fig:02-09c}]{\includegraphics[height=5cm]{Fig09c}}\hfill
\subfloat[\label{fig:02-09a}]{\includegraphics[height=5cm]{Fig09a}}\hfill
\subfloat[\label{fig:02-09b}]{\includegraphics[height=5cm]{Fig09b}}
\end{centering}
\caption{Exemples de structures : \ref{fig:02-09c}  1D (Olivine \ce{LiFePO4}) ; \ref{fig:02-09a} 2D (Lamellaire \ce{LiCoO2}) ; \ref{fig:02-09b} 3D (Spinelle \ce{LiMn2O4})}
\label{fig:09}
\end{figure}

\definition{Les structures moléculaires d'insertion se distinguent par le \blancbis{nombre de directions} orthogonales par lesquelles un ion inséré peut \blancbis{diffuser}. Ce critère permet de les \blancbis{catégoriser}.}

Comme représenté à la Fig. \ref{fig:09}, la diffusion des ions lithium peut avoir lieu :
\begin{itemize}
\item dans la direction $\vec{b}$ dans l'olivine ;
\item dans les directions $(\vec{a}, \vec{b})$ dans les matériaux lamellaires ;
\item dans les trois directions dans la spinelle.
\end{itemize} 

Dans le domaine des batteries, on fait une distinction physique entre la \blancbis{charpente moléculaire}, laquelle existe toujours que l'on charge/décharge et les \blancbis{sites} accueillant le \blancbis{lithium} (qui peut être plus ou moins présent en fonction du SoC).

En fonction du type de structure rencontré, plusieurs termes de vocabulaire sont usités par les scientifiques pour décrire les structures : 

\definition{Les termes fréquemment rencontrés sont les suivants : 

\begin{center}
\begin{tabular}{cccc}
\toprule
 & Termes anglais & Termes français & Exemple de structures \\
\midrule
\multirow{4}{*}{Charpente} & Sheets & Feuilles & Graphène \\
 & Stacks & Pile de feuilles & Graphite \\
 & Slabs & Plaques & Lamellaires \\
 & Pillars & Pilliers & Olivines \\
\midrule
\multirow{4}{*}{Insertion} & Galeries & Galeries & Lamellaires \\
 & Tunnels & Tunnels & Olivine ou Spinelle \\
 & Cavities & Cavités & - \\
 & Windows & Fenêtres & - \\
\bottomrule
\end{tabular}
\captionof{table}{Termes fréquemment rencontrés pour décrire les structures d'insertion/intercalation du lithium}
\end{center}}

\subsection{Environnement chimique du lithium inséré : influence sur la f.e.m.}

Lorsque l'on étudie les performances de matériaux de différentes dimensions, on se rend compte qu'en théorie, tous les matériaux d'une même classe disposent d'une  capacité théorique proche\mySideNote{éléments de transition proches dans le tableau périodique}.

\begin{center}
\begin{tabular}{cccc}
\toprule
\multirow{2}{*}{Dimension} & \multirow{2}{*}{Composé} & Capacité spécifique  & Potentiel rédox moyen \\ 
& & [\si{\milli\ampere\hour\per\gram}] & \si{\volt}/\ce{Li^+/Li}  \\
\midrule
\multirow{2}{*}{1D} & \ce{LiFePO4} & 170 (160) & 3.45 \\
 & \ce{LiFe_{0.5}Mn_{0.5}PO4} & 170 (160) & 3.4/4.1 \\
\midrule
\multirow{2}{*}{2D} & \ce{LiCoO2} & 272 (140) & 4.2 \\
& \ce{LiNi_{0.33}Mn_{0.33}Co_{0.33}O2} & 272 (200) & 4.0 \\
\midrule
\multirow{2}{*}{3D} & \ce{LiMn2O4} & 148 (120) & 4.1 \\
 & \ce{LiMn_{1.5}Ni_{0.5}O4} & 148 (120) & 4.7 \\
\bottomrule
\end{tabular}
\captionof{table}{Capacités spécifiques et potentiels rédox de quelques structures cristallographiques lithiées.}
\end{center} %Inorganics 2014, 2

En effet, si deux matériaux lithiés ont la même \blancbis{charpente moléculaire}\mySideNote{donc ils appartiennent au même groupe d'espace}, on peut s'attendre théoriquement à ce qu'ils puissent \blancbis{accueillir autant} d'ions lithium. Cependant, en pratique, ceci n'est pas forcément vérifié expérimentalement à cause d'éventuelles différences de stabilité de ces deux structures vis-à-vis de la présence d'ions lithium insérés.

Le \blancbis{potentiel rédox}, lui, est très dépendant de la \blancbis{nature} des atomes, particulièrement des \blancbis{cations}\mySideNote{mais aussi indirectement des anions comme nous le verrons au dernier chapitre} de la charpente qui changent l'environnement chimique des atomes de lithium\mySideNote{comme l'hydratation des ions lithium impacte les performances rédox par rapport à la phase gaz}. Cet environnement chimique va jouer un rôle crucial dans la \blancbis{stabilité} des phases cristallines.

\definition{La capacité spécifique est liée à la structure, tandis que la force électromotrice est liée à la force des liaisons chimiques mises en jeu dans les structures.}

%\section{Diffusion du Li au sein d'un hôte et stabilité thermodynamique}

\subsection{Critères de stabilité des sites et mécanisme de diffusion}

Lorsque le lithium va s'insérer dans un matériau, il peut apparaître de multiples phases stables thermodynamiquement\mySideNote{Nous nous contenterons d'interpréter des stabilités d'ordre thermodynamique. Bien entendu, la vitesse d'insertion trop rapide des ions lithium peut conduire à la formation de phases instables thermodynamiquement, mais qui se forment plus rapidement}.

Pour les interpréter, il faut comparer les phases thermodynamiquement stable : 

\begin{center}
\includegraphics[width=0.9\textwidth]{Fig30}
\captionof{figure}{Réaction "fictive" de mélange d'une phase contenant que du lithium à une phase ne contenant pas de lithium}
\end{center}

Dans cette réaction, l'enthalpie de mélange $\Delta G_m$ associée à la différence d'enthalpie entre la phase sans lithium et la phase avec lithium s'écrit à une température et une pression donnée :

\begin{align*}
\Delta G_m (x) = G (x) - \left[ x G(x = 1) + (1-x) G(x=0) \right]
\end{align*}

avec $x$ la fraction de lithium inséré dans le matériau. Le signe de $\Delta G_m$ à une fraction $x$ de lithium nous informe sur le gain de stabilité apporté lorsque l'on ajoute du lithium dans la structure. Or, cette enthalpie de mélange dépend de deux facteurs : 

\begin{align*}
\Delta G_m (x) = \Delta H_m (x) - T \Delta S_m (x)
\end{align*}

avec $\Delta H_m$ représentant le terme enthalpique et $\Delta S_m$ représentant le terme entropique.

\`{A} un courant modéré, si le lithium ne perturbe pas\mySideNote{n'interagit pas} la structure, l'\blancbis{entropie} permettra toujours au lithium de s'insérer/de se solubiliser dans une structure présentant des sites d'insertion.

Cependant, en général, l'insertion ou la desinsertion de lithium peut créer une instabilité \blancbis{enthalpique} de la structure : il se passe alors un phénomène démixion, c'est-à-dire l'apparition de plusieurs phases thermodynamiquement stables, l'une enrichie en lithium et l'une appauvrie en lithium.

\begin{figure}[H]
\begin{centering}
\subfloat[\label{fig:02-14a}]{\includegraphics[height=5cm]{Fig14a}}\hfill
\subfloat[\label{fig:02-14b}]{\includegraphics[height=5cm]{Fig14b}}\hfill
\subfloat[\label{fig:02-14c}]{\includegraphics[height=5cm]{Fig14c}}
\caption{\ref{fig:02-14a} Diffusion aléatoire de l'ion inséré sans modification de structure ; \ref{fig:02-14b} Diffusion aléatoire de l'ion intercalé avec apparition de deux phases distinctes (avec et sans lithium) ; \ref{fig:02-09b} Diffusion ordonnée de l'ion intercalé avec apparition de deux phases distinctes (avec et sans lithium)}
\end{centering}
\end{figure}

Le mécanisme de diffusion est largement dépendant de l'effet enthalpique du lithium sur la structure :

\begin{center}
\begin{minipage}{0.40\textwidth}
\begin{tabular}{cccc}
\toprule
& \ref{fig:02-14a} & \ref{fig:02-14b} & \ref{fig:02-14c} \\
\midrule
$\Delta H_m$ & $\sim 0$ & $\gg 0$ & $\ll 0$ \\
$-T \Delta S_m$ & $<0$ & $<0$ & $<0$ \\
$\Delta G_m$ & $< 0$ & $> 0$ & $\ll 0$ \\
\bottomrule
\end{tabular}
\end{minipage}\begin{minipage}{0.56\textwidth}
\captionof{table}{Types de mécanismes attendus en fonction de la stabilité de la structure lors de l'insertion du lithium}
\end{minipage}
\end{center}

Expérimentalement, on peut retranscrire l'effet de la température et de la fraction de lithium à l'aide d'une courbe de solubilité, appelée courbe du \blancbis{solvus} : 

\begin{center}
\begin{minipage}{0.48\textwidth}
\centering
\includegraphics[width=0.5\textwidth]{Fig31}
\end{minipage}\begin{minipage}{0.48\textwidth}
\captionof{figure}{Allure d'une courbe de solvus en fonction de la fraction en lithium dans la structure.}
\end{minipage}
\end{center}

En général dans les solides, la température favorise l'insertion du lithium dans une structure. Au-dessus d'une température critique, appelée \blancbis{température de consolution}, le système est monophasique pour toute composition en lithium. On se retrouve alors dans la situation de la Fig. \ref{fig:02-14a}.

Sinon, le système peut avoir une place de composition pour laquelle il existe de manière biphasique. On se retrouve alors dans la situation de la Fig. \ref{fig:02-14b}.

\definition{
Les mécanismes d'insertion du lithium sont largement dépendants de la stabilité de la structure lors de l'insertion. Le terme \blancbis{enthalpique} joue un rôle crucial dans les mécanismes mis en jeu et dépend exclusivement :
\begin{itemize}
\item la stabilité de la charpente moléculaire en présence ou absence de lithium ;
\item la stabilité du lithium dans les sites d'insertions.
\end{itemize}
}


%Ainsi, on peut dresser un bilan du type de mécanisme que l'on peut observer pour le lithium en fonction de deux considérations précédentes (stabilité des sites/stabilité de la charpente) : 

%\begin{center}
%\begin{tabular}{cccc}
%\toprule
%& \ref{fig:02-14a} & \ref{fig:02-14b} & \ref{fig:02-14c} \\
%\midrule
%Charpente initiale stable avec Li & \faCheck & \faRemove & \faRemove \\
%Sites préférentiels & \faRemove & \faRemove & \faCheck \\
%\bottomrule
%\end{tabular}
%\captionof{table}{Types de mécanisme attendu en fonction de la stabilité du lithium dans les différents sites d'insertion et la stabilité de la charpente moléculaire en présence d'ions lithiums}
%\end{center}
%Charpente alternative stable avec Li & \faCheck ou \faRemove & \faCheck & \faCheck \\

%Dans l'ensemble, cela détermine les différents modes de diffusion. Lors de l'insertion du lithium, si celui-ci est toléré dans la structure hôte mais que tous les sites sont équivalens, on se retrouve dans la situation \ref{fig:14a}. Dans le cas où une transition de phase permet à la structure d'accueillir le lithium, mais que tous


\subsection{Conséquences sur la force électromotrice}

Nous avons vu que les courbes de charge/décharge contiennent une information thermodynamique via la force électromotrice. Dans le système :
\begin{equation*}
\oplus ~ \ce{MO} (s) | \ce{Li^+} (\text{electrolyte}) | \ce{Li(s)} ~ \ominus
\end{equation*} 

la force électromotrice nous donne le potentiel d'électrode au pôle $\oplus$, laquelle sera liée à la nature des phases présentent au sein du matériau. On a du point de vu des phases, trois cas possibles illustrés à la section précédente : 
\begin{enumerate}
\item formation d'une seule phase homogène où le lithum diffuse aléatoirement au sein du matériau (c.-à-d. une phase \ce{Li_x MO}) (cf. Fig. \ref{fig:02-14a}) ; 
\item formation de deux phases, non miscibles, où le lithium diffuse aléatoirement au sein du matériau (c-à-d. formation d'une phase \ce{Li MO} et une phase \ce{MO}) (cf. Fig. \ref{fig:02-14b}) ;
\item formation de deux phases, partiellement miscibles, où le lithium diffusion aléatoire au sein du matériau (c-à-d. formation d'une phase enrichie en lithium \ce{Li_x MO} et une phase appauvrie \ce{Li_y MO} où $x \gg y$) (cf. Fig. \ref{fig:02-14c}).
\end{enumerate}

\begin{figure}[H]
\begin{centering}
\subfloat[\label{fig:23a}]{\includegraphics[height=4.5cm]{Fig23a}}\hfill
\subfloat[\label{fig:23b}]{\includegraphics[height=4.5cm]{Fig23b}}
\end{centering}
\caption{\'{E}volution du potentiel d'électrode $\oplus$ en fonction de la composition $x$ de l'électrode en Li dans le cas \ref{fig:23a} monophasé et dans le cas \ref{fig:23b} biphasé}
\end{figure}

\definition{Les courbes de charge/décharge sont donc capable de nous renseigner sur les phases présentent, ainsi que les transitions de phases à une température donnée.}

\section{Stabilité de la structure hôte et transition de phase}

\subsection{Stabilisation d'une charpente moléculaire}

Afin de réflechir à la \blancbis{stabilité} d'une \blancbis{charpente moléculaire}, nous allons poser une structure \blancbis{lamellaire} \ce{MX2} où M est le \blancbis{métal} et X un \blancbis{anion divalent} (oxygène, soufre, sélénium).%adoptant une structure lamellaire.

%Comme vu dans l'étude de la stabilité des sites du lithium, la répulsion coulombienne entre les feuillets est un critère de stabilité important. La stabilisation d'une charpente lamellaire peut se faire de deux manières :
Dans les structures lamellaires, la stabilité d'une charpente lamellaire est liée à la répulsion coulombienne entre les plaques, qui peut être diminuée en :
\begin{enumerate}
\item réduisant la charge partielle des anions \ce{X} au sein de la structure,
\item écrantant la charge partielle des anions \ce{X}.
\end{enumerate}

Pour cela, diverses stratégies sont envisageables par le chimiste : 
\begin{enumerate}
\item utiliser des anions plus gros (\emph{par ex.} soufre, sélénium), permettant de rendre plus \blancbis{diffuse} la charge $\ominus$ des anions (cf. \ref{fig:16a}) ;
\item utiliser des \blancbis{métaux de transition} avec un grand degré d'oxydation\mySideNote{ou nombre d'oxydation} qui augmentera la \blancbis{donation d'électrons} de l'anion $X$ vers le métal $M$ (cf. \ref{fig:16b}) ;
\item utiliser des structures où il existe une fraction non négligeable de \blancbis{cations} dans les galeries pour stabiliser la charpente (cf. \ref{fig:16c}).
\end{enumerate}

\begin{figure}[H]
\begin{centering}
\subfloat[\label{fig:16a}]{\includegraphics[height=3cm]{Fig16a}}\hfill
\subfloat[\label{fig:16b}]{\includegraphics[height=3cm]{Fig16b}}\hfill
\subfloat[\label{fig:16c}]{\includegraphics[height=3cm]{Fig16c}}
\end{centering}
\caption{\ref{fig:16a} Structures soufrés ; \ref{fig:16b} Structures à haut d.o. ; \ref{fig:16c} Cations \ce{Li+} ou \ce{H+} présents pour stabiliser la structure.}
\end{figure}

Dans le cas des autres structures : olivine et spinelle, l'absence de galeries implique beaucoup moins de modification structurales lors de l'ajout ou du départ du lithium.

\subsubsection{Paramètres influençant la stabilité du lithium au sein de la structure hôte}

Nous allons nous appuyer sur les structures lamellaires pour définir quelques règles de bon sens justifiant la stabilité des cations lithium au sein d'une de ces structures.

\begin{minipage}{0.45\textwidth}
\centering
\includegraphics[width=6cm]{Fig10}
\end{minipage}\hfill\begin{minipage}{0.45\textwidth}
\captionof{figure}{Trois polyèdres de coordination couramment rencontrés dans les structures lithiées}
\label{fig:10}
\end{minipage}

\definition{L'environnement de coordination du lithium est défini par le polyèdre formé par les anions qui lui sont les plus proches voisins.}

Comme illustré à la Fig. \ref{fig:10}, dans les structures lamellaires, on trouve trois types de polyèdres de coordination\mySideNote{Coordinence d'un atome \ce{A} est le nombre d'atomes qui sont proches voisins de lui.} : l'octaèdre [6], le prisme [6] et le tétraèdre [4].

Au sein des structures lamellaires, les scientifiques ont par exemple étudié les structures lithiées \ce{LiCoO2} et \ce{LiTiS2}.

\begin{minipage}{0.65\textwidth}
\centering
\includegraphics[width=10cm]{Fig11}
\end{minipage}\hfill\begin{minipage}{0.25\textwidth}
\captionof{figure}{Trois empilements rencontrés dans les structures \ce{LiCoO2} et \ce{LiTiS2}}
\label{fig:11}
\end{minipage}

\`{A} la Fig. \ref{fig:11}, trois structures O1, O3 et P3\mySideNote{Cette notation se lit comme suit : 
\begin{enumerate}
\item une lettre pour le polyèdre formé par l'atome inséré ;
\item un nombre pour le nombre minimal de plaques à inclure dans la maille élémentaire.
\end{enumerate}
} communes sont présentées. Ces structures correspondent à des \blancbis{empilements compacts} d'atomes d'oxygène différent (resp. $BABABA\ldots$, $BACBACBA\ldots$, $BAACCB\ldots$), séparées par des galeries. Ces trois \blancbis{variétés allotropiques} de \ce{LiCoO2} sont structuralement proches, puisque l'on peut passer de l'une à l'autre par un \blancbis{simple glissement} relatif des \blancbis{plaques}. On peut donc s'attendre à rencontrer ces trois structures pour tout état de charge et puisqu'elles sont proches, qu'elles puissent coexister simultanément\mySideNote{Comme le dit l'adage britannique, "like dissolves like"}.

\begin{minipage}{0.45\textwidth}
\centering
\includegraphics[width=6cm]{Fig13}
\end{minipage}\hfill\begin{minipage}{0.45\textwidth}
\captionof{figure}{Distance des ions lithium présent au sein d'une même galerie [représentation selon une projection dans le plan d'une galerie]}
\label{fig:13}
\end{minipage}

Dans ces structures, deux environnements de coordination du lithium sont présents : octaédrique et prismatique. Si on regarde comment les ions lithium sont organisés dans le plan de la galerie (voir Fig. \ref{fig:13}), on remarque un réseau en nid d'abeilles (resp. triangulaire) pour une coordination prismatique (resp. octaédrique). La structure P3 permet donc d'augmenter la distance cation-cation au sein de la galerie.

\begin{minipage}{0.55\textwidth}
\centering
\includegraphics[width=8cm]{Fig12}
\end{minipage}\hfill\begin{minipage}{0.35\textwidth}
\captionof{figure}{Comparaison du rapport entre le rayon du cation alcalin et du rayon des anions au sein d'un octaèdre et d'un tétraèdre.}
\end{minipage}

Cependant, la structure P3 n'est pas celle observée à l'état \SI{100}{\percent} chargé, ce n'est \emph{a priori} pas la plus stable. En effet, si on considère qu'il y a contact anion-anion et entre le lithium-anion, on peut alors calculer l'espacement entre deux plaques\mySideNote{slabs}. On observe que celle-ci est plus courte pour un prisme que pour un octaèdre. Cet effet déstabilisant est d'autant plus grand que le cation est petit\mySideNote{or le lithium est le plus "petit" alcalin.}.

Enfin, un autre aspect qui peut-être abordé est la connectivité entre les polyèdres de lithium et des cations de la charpente. On peut compter le nombre de faces mises en jeu : O3 < O1 < P3\mySideNote{avec "<" qui signifie "a moins de face commune entre polyèdres que"}.

\definition{
Lorsque les ions lithium s'insérent dans une structure, ils le font en général en obéissant à quelques règles de bon sens :

\begin{enumerate}
\item minimisant la répulsion électrostatique entre les cations lithium insérés ;
\item minimisant la répulsion électrostatique entre les cations lithium et ceux de la charpente ; 
\item minimisant la répulsion entre anions coordonnés à l'ion lithium (donc entre les plaques lorsque la structure est lamellaire).
\end{enumerate}}


\begin{minipage}{0.55\textwidth}
\centering
\includegraphics[width=8cm]{Fig15}
\end{minipage}\hfill\begin{minipage}{0.35\textwidth}
\captionof{figure}{Courbe de charge de \ce{Li_xCoO2} à courant constant valant $C/24$ à température ambiante.}
[H1-3 est une phase intermédiaire où coexistent les phases O1 et O3]
\end{minipage}%WARNING

Dans \ce{Li_x CoO2}, il a été observé que : 
\begin{itemize}
\item O1 est la phase la plus stable lorsqu'il n'y a pas de lithium ($x \to 0$);
\item O3 est la phase la plus stable lorsqu'il y a du lithium ($x \to 1$) ;
\item P3 n'est sont pas observé mais cette structure existe pour de plus gros cations ($\ce{Na^+}, \ce{K^+}\ldots$).
\end{itemize}

\subsection{Interprétation de la stabilité des liaisons chimiques}

Il existe des règles permettant de rationnaliser la stabilité des structures cristallographiques : que ça soit au niveau de la charpente moléculaire ou au niveau des ions insérés. Cependant, ces règles dépendent de la nature de la liaison chimique : covalente ou ionique pour les oxudes.

\begin{center}
\begin{minipage}{0.48\textwidth}
\centering
\includegraphics[width=0.8\textwidth]{Fig29}
\end{minipage}\hfill\begin{minipage}{0.48\textwidth}
\captionof{figure}{Triangle de Katelaar}
\label{fig:29}
\end{minipage}
\end{center}

Comme le représente le triangle de Katelaar à la Fig. \ref{fig:29}, l'ionicité de la liaison dépend de la différence d'électronégativité entre les atomes. Plus cette différence est grande, plus les liaisons chimiques sont proches d'un modèle purement ionique. Dans le cas contraire, la liaison possède un caractère covalent, le pôle métallique n'étant pas accessible pour un oxyde ou dérivé.



\subsubsection{Caractère ionique : critères structuraux}
%Solid State Chemistry and its Applications
Des règles \blancbis{structurales} ont été pausé par \blancbis{Pauling} en 1929 pour rationaliser la stabilité d'une structure cristallographique. Ces règles sont au nombre de 5\mySideNote{Ces règles ne sont pas absolues, il n'existe que très peu de structures qui respectent l'ensemble des règles. Toutefois, elles restent un outil rapide pour rationaliser le comportement des structures.} :  

\definition{
\blancbis{Règle du rapport des rayons} 

La distance entre un cation et un anion est déterminée par la somme des rayons \blancbis{ioniques}, soit $r_\oplus + r_\ominus$. La coordinence du polyèdre de coordination dépend du rapport des rayons : 
\begin{align*}
\frac{r_\oplus}{r_\ominus}
\end{align*}}

\begin{center}
\includegraphics[width=10cm]{Fig17}
\captionof{figure}{\'{E}volution attendue de la coordination en fonction du rapport des rayons}
\end{center}

Cette règle, déjà vue en première année, vous permet de "prédire" des espèces en fonction du rapport de leurs rayons ioniques. Pour cela, il faut établir deux relations permettant de se ramener aux paramètres de maille.

\definition{\blancbis{Règle de la valence} 

Dans un polyèdre de coordination, la liaison chimique anion-cation est d'autant plus forte que le rapport :

\begin{align*}
\frac{\text{charge du cation}}{\text{coordination du cation}}
\end{align*} 

est grand.}

\begin{figure}[H]
\begin{centering}
\subfloat[\label{fig:02-18a}]{\includegraphics[width=3cm]{Fig18a}}\hfill
\subfloat[\label{fig:02-18c}]{\includegraphics[width=3cm]{Fig18c}}\hfill
\subfloat[\label{fig:02-18b}]{\includegraphics[width=3cm]{Fig18b}}
\end{centering}
\caption{Classement de 3 coordinations par ordre croissant de stabilité en allant de gauche à droite.}
\end{figure}

Cette règle vous permet de "prédire" la force des liaisons ioniques dans un matériau. Elle suppose qu'entre deux structures cristallographiques équivalentes, les effets de dilatation ou contraction de la structure sont négligeables.

\definition{\blancbis{Règle de connectivité des polyèdres}

La présence d'arêtes partagées, et en particulier de faces partagées, dans une structure coordonnée diminue sa stabilité. Cet effet est d'autant plus important pour les cations ont une \blancbis{grande valence} (un grand d.o.) et un petit \blancbis{nombre de coordinations}, et il est particulièrement important lorsque le rapport des rayons approche \blancbis{la limite inférieure} de stabilité du polyèdre.}

Cette règle est assez importante, et permet de facilement identifier la stabilité des polyèdres entre eux en fonction de leur \blancbis{connectivité}.

\definition{\blancbis{Règle des "mélanges" de cations} Dans un cristal contenant différents cations, ceux qui ont une grande valence et un petit nombre de coordinations ont tendance à ne pas avoir de connexions les uns les autres.}

\begin{minipage}[]{0.5\textwidth}
\centering
\includegraphics[width=5cm]{Fig19}
\end{minipage}\hfill\begin{minipage}[]{0.46\textwidth}
\centering
\captionof{figure}{Projection de la structure d'une olivine Forstérite \ce{Mg2SiO4}}
\end{minipage}

Par exemple, lorsque l'on regarde la structure d'une olivine dans un plan, on constate que les polyèdres de silicium ne sont jamais directement reliés entre eux.

\definition{\blancbis{Règle de parcimonie} Le nombre de polyèdres différents dans un cristal tend à être faible.}

En général, on trouve 1 ou 2 polyèdres différents dans un cristal. Ceci se comprend bien si on considère qu'un cation va privilégier un environnement chimique\mySideNote{un polyèdre de coordination} donné.

On remarquera que l'on a utilisé des idées équivalentes pour interpréter la stabilité du lithium, peut-être est-ce dû au caractère assez ionique de la liaison \ce{Li-O} ?%WARNING

\subsubsection{Caractère ionique : critères énergétiques}

La stabilité d'une structure à fort caractère ionique peut être expliquée à l'aide du formalisme de la \blancbis{théorie du champ cristallin}\mySideNote{cf annexe \ref{chp:cristal_field}, vu en BUT1}. Cette théorie peut permettre d'interpréter la stabilité des \blancbis{oxydes métalliques} utilisés dans les batteries, comme par exemple : 
\begin{itemize}
\item pour prédire si une structure est \blancbis{spinelle} ou \blancbis{anti-spinelle} dans les métaux de transition ;
\item pour prédire si le \blancbis{polyèdre de coordination} est \blancbis{octaédrique} ou \blancbis{tétraédrique}.
\end{itemize}

\begin{center}
\includegraphics[width=0.9\textwidth]{Fig20}
\captionof{figure}{Diagramme d'un atome $d^6$ (\emph{par ex.} le Co(III)) en configuration champ fort/bas spin}
\end{center}

Par exemple, dans le cas de la structure \ce{Li_xCoO2}, si on analyse les structures limites\mySideNote{correspondant à l'état $SoC=$\SI{100}{\percent} et $SoC=$\SI{0}{\percent})}, on est en capacité de calculer la différence d'énergie entre ces deux structures, et donc d'estimer la stabilité relative selon le champ cristallin. 

\begin{center}
\begin{adjustbox}{max width=\textwidth}
\begin{tabular}{ccccccc}
\toprule
Nom & $no(\ce{Co})$ & Configuration & Type de champ & $ESCC(Oh)$ & $ESCC(Td)$ & $\Delta ESCC$ \\
\midrule
\ce{LiCoO2} & III & $d^6$ & Bas spin/Champ fort & $-2.4 ~ \Delta_O$ & $-0.6 ~ \Delta_T$ & $-2.13 ~ \Delta_O$ \\
\ce{CoO2} & IV & $d^5$ & Bas spin/Champ fort & $-2.0 ~ \Delta_O$ & $0$ & $-2.00 ~ \Delta_O$ \\
\bottomrule
\end{tabular}
\end{adjustbox}
\captionof{figure}{Calcul de la différence de stabilité entre des structures tétaédrique (Td) et octaédriques (Oh) pour le Co(III/IV)}
[l'énergie d'appariement des électrons est négligée, $\Delta ESCC$ est défini comme $ESCC(Oh) - ESCC (Td)$]
\end{center}

On montre ainsi que dans les deux structures limites, le \ce{Co} est \emph{a priori} dans un environnement octaédrique.

\begin{center}
\includegraphics[width=10cm]{Fig21}
\captionof{figure}{Diagramme d'un atome $d^4$ (\emph{par ex.} le Mn(III)) en configuration champ faible/haut spin}
\end{center}

La théorie du champ cristallin est aussi parfois un bon moyen d'interpréter les déformations ayant lieu au sein des polyèdres de coordination. Par exemple, dans le cas d'un octaèdre, il est courant que les configurations $d^4$ (champ faible), $d^7$ (champ fort) et $d^9$ se stabilisent en ayant un polyèdre irrégulier (octaèdre étiré). Cet effet est communément appelé effet \blancbis{Jahn-Teller}.

\definition{La théorie du champ cristallin est une méthode appropriée pour interpréter énergétiquement la stabilité des polyèdres de coordination présents dans les oxydes ou dérivés utilisés dans les batteries. Cependant, la nature \blancbis{ionique} des liaisons chimiques sous-tend son utilisation.}

\subsubsection{Caractère covalent}

Au sein des oxydes métalliques, les oxygènes forment des \blancbis{liaisons chimiques} avec les métaux, souvent des métaux de transition.

L'écart d'électronégativité entre les anions et les métaux de transition est souvent insuffisante pour conduire à la formation de liaisons chimiques purement ionique, comme représenté à la Fig. \ref{fig:02-07c}.\mySideNote{Au contraire, on s'attend à ce que la liaison \ce{Li-O} soit proche de ce qui est représenté à la Fig. \ref{fig:02-07c}.}

\begin{figure}[H]
\begin{centering}
\subfloat[\label{fig:02-07a}]{\includegraphics[height=5cm]{Fig07a}}\hfill
\subfloat[\label{fig:02-07b}]{\includegraphics[height=5cm]{Fig07b}}\hfill
\subfloat[\label{fig:02-07c}]{\includegraphics[height=5cm]{Fig07c}}
\end{centering}
\caption{Diagramme d'orbitales moléculaires entre une entité atomique \ce{A} et une entité atomique \ce{B} attendu pour une liaison \ref{fig:02-07a} purement covalente, \ref{fig:02-07b} iono-covalente et \ref{fig:02-07c} purement ionique}
[$\Delta \chi = \chi_B - \chi_A$ montre que l'évolution du diagramme orbitalaire dépend de la différence d'électronégativité des entités atomiques, avec ici \ce{B} plus électronégatif que \ce{A}]
\end{figure}

Par conséquent, comme représenté à la Fig{fig:02-07b}, les liaisons formées sont généralement des liaisons \blancbis{iono-covalentes} où l'utilisation de modèles prévu pour les liaisons ioniques (règles de Pauling, champ cristallin) ne fonctionnent pas ou ne permettent pas de justifier tous les résultats expérimentaux.

Ainsi, au second ordre, on pourra être emmené à interpréter des \blancbis{diagrammes d'orbitales} pour interpréter les déviations entre l'expérience et le modèle ionique.\mySideNote{Dans ce cours, on se limitera à la construction d'interactions orbitalaires simples entre deux orbitales atomiques$\ldots$} 

\definition{Pour interpréter la stabilité des liaisons possédant un caractère \blancbis{covalent}, à l'aide de diagrammes orbitalaires, on aura besoin :
\begin{enumerate}
\item de représenter les orbitales $s$, $p$ et $d$ ;
\item de connaître l'influence de l'électronégativité sur la position des orbitales $s$, $p$ et $d$ d'un atome ;
\item de construire des interactions à deux orbitales ;
\item d'interpréter la distribution des électrons dans les orbitales liantes et anti-liantes résultant d'une interaction à deux orbitales.
\end{enumerate}
}

\begin{comment}
%\subsection{Comportement des courbes de titrage coulométriques}

%A voir !

Sites Td versus sites Oh

Structure cristallographique : site pour accueillir un ion Li

Formes élémentaires : Td, Oh


OD : 
1D : NaCoSO4F.2H2O, MX3, AFeS2, A Mo3X3, nanotubes de C
2D : CoO2, graphite, dichalcogénures, oxyhalogénures
3D : LiFeSO4F, spinelles, phases de chevrel, péroxskites, rutiles

Li+ : sites d'accueil (Td, trigonal prismatique, Oh) , réaction d'inversion : 
Plus la répulsion est élevée : sites octaedriques plutôt que trigonal prismatique.
Ex : TiS2, NbS2 (charge soufre-séleniures diffuse)



V2O5, LiCoO2, Ni(OH)2
-> augmenter le no du métal : V(+V)
-> intercalation du lithium (synthèse in situ avec du lthium)
-> proton

%Structure de LiMn2O4
%https://www.youtube.com/watch?v=2_IILk7jzkQ&ab_channel=MaterialsScienceEducation
%Phase transtion Cuvbic/Tetragonal


Lien : Maille élémentaire -> Cristal

Electrochimie : Dualité ion - électron 
\end{comment}

\begin{comment}
\section{Relation entre structure cristallographique et transfert des électrons}

Lorsque l'on étudie les réactions rédox, il va de soi qu'il est nécessaire de parler de structure électronique. Or, au sein de molécules, la structure électronique peut être décrite à l'aide de la théorie des \blancbis{orbitales moléculaires}\mySideNote{abordée en BUT1}, résultats pouvant être qualitativement décrits, feuille et crayon en main, à l'aide de la méthode des \blancbis{combinaisons linéaires d'orbitales atomiques}.

\section{Formalisme simplifié de la théorie des bandes}

Dans le cas des solides cristallins, l'approche est bien plus complexe. En effet, vous avez vu qu'un solide peut être décrit comme un \blancbis{enchaînement infini} de \blancbis{motifs cristallins} sur le réseau périodique. De cela, il en résulte qu'un solide peut-être vu comme une molécule géante composée d'une infinité d'atomes, espacés de manière constante et régulière sur le réseau.

\subsection{Construction du modèle}

Nous n'allons pas chercher à construire un modèle de bande parfait, nous nous limiterons aux notions importantes pour comprendre les comportements électroniques dans les solides cristallins, et plus précisément les \blancbis{oxydes métalliques} ou autres dérivés analogues.

Prenons un cas simple : une anode en lithium métallique. Pour simplifier, nous considérerons un réseau 1D où les sites du réseau sont espacés d'un pas $a$ constant\mySideNote{la distance \ce{Li-Li}}.

\begin{center}
\includegraphics[height=5cm]{Fig04}
\captionof{figure}{Construction d'une bande à l'aide d'un réseau unidimentionnel présentant une périodicité $a$ et une longueur infinie.}
\end{center}

La structure électronique de l'atome de lithium nous indique qu'il ne possède qu'une orbitale de valence, simplement occupée\mySideNote{l'orbitale $2s$}. Lorsque l'on ajoute progressivement des atomes de lithium sur le réseau, on réalise des combinaisons linéaires d'orbitales atomiques, on a alors : 

\begin{itemize}
\item formation de niveaux \blancbis{liants} et \blancbis{antiliants} ;
\item rapprochement des niveaux au fur et à mesure de l'ajout des atomes.
\end{itemize}

Lorsque l'on atteind une infinité d'atome\mySideNote{Cela n'a pas de sens physique, pour un physico-chimiste, on peut appréciablement considérer que le nombre d'atomes dans une molécule tend vers $\infty$ lorsque le nombre d'atomes mis en jeu en commensurable au nombre d'Avogadro $N_A$.}, les niveaux sont infiniment proche\mySideNote{Cela signifie physiquement que la différence d'énergie entre deux niveaux est plus petite que l'énergie thermique (de l'ordre de $\sim RT$).}. On a alors formation d'une \blancbis{bande continue d'énergie}.

\subsection{Modèle des bandes au c{\oe}ur des transferts rédox}
Le transfert d'électron implique le plus souvent des matériaux d'électrodes, dont la structure électronique est décrite par une ou plusieurs bandes. On considère initialement une électrode $\ominus$ composée de lithium métallique.

\begin{minipage}{0.5\textwidth}
\includegraphics[height=4.5cm]{Fig05b}
\end{minipage} \hfill \begin{minipage}{0.45\textwidth}
\captionof{figure}{Bande simple attendue pour une électrode en Li au pôle $\ominus$}
\label{fig:05b}
\end{minipage}

Lorsque deux atomes de \ce{Li} interagissent, on sait que ce dimère est de configuration :
\begin{align*}
(\sigma)^2 (\sigma^{\star})^0
\end{align*}

où $\sigma$ désigne le niveau liant (plein) et $\sigma^\star$ désigne le niveau anti-liant (vide). Il existe donc 4 spin-orbitales\mySideNote{On considère une orbitale comme une fonction électronique décrite par le triplet $(n,l,m)$. On appelle spin-orbitale, la fonction électrique à laqulle on ajoute une composante décrivant l'état de spin, soit un quadruplet $(n,l,m,m_s)$} possibles pour les électrons, et seulement la moitié est occupée. Si on étend ceci à $N$ atomes, on s'attend donc à avoir $N$ orbitales dont $N/2$ sont occupées.

Un peut alors poser trois notions élémentaires pour décrire une bande : le niveau de Fermi, la bande de conduction et la bande de valence.

\definition{
\blancbis{Niveau de Fermi} Dans un solide à \SI{0}{\kelvin}, le niveau de Fermi désigne la limite en énergie entre le dernier niveau occupé et le premier niveau vide. Il est noté $\epsilon_F$ et est donné le plus souvent en \si{\electronvolt}.

~

\blancbis{Bandes de conduction} La(es) bande(s) de conduction désigne(nt) l'ensemble des orbitales innocupées. Si on pose $\epsilon_c$, l'énergie du niveau le plus bas, la bande de conduction est définie sur l'intervalle $[\epsilon_c, + \infty[$.

~

\blancbis{Bandes de valence} La(es) bande(s) de valence désigne(nt) l'ensemble des orbitales occupées. Si on pose $\epsilon_v$, l'énergie du niveau le plus haut, la bande de valence est définie sur l'intervalle $]- \infty,\epsilon_v]$.

~

Le niveau de Fermi peut se définir relativement aux bandes, et équivaut à :

\begin{align*}
\epsilon_F \approx \frac{\epsilon_c + \epsilon_v}{2}
\end{align*}
} 

Dans le cas du lithium métallique, on a une correspondance directe entre niveaux liants (resp. antiliants) et la bande de valence (resp. de conduction). Lorsque l'on oxyde le lithium, on a \blancbis{perte d'électrons} dans les niveaux liants de la bande de valence. Lorsque l'on réduit le lithium, on a \blancbis{gain d'électrons} dans les niveaux antiliants de la bande de conduction. Dans les deux situations, l'oxydation ou la réduction a pour effet de \blancbis{destabiliser} une structure moléculaire \blancbis{stable}.

Du point de vu énergétique, le niveau de Fermi est une sonde du \blancbis{potentiel} de l'électron au sein d'un solide comme illustré aux Fig. \ref{fig:05a} et Fig. \ref{fig:05c}.

\begin{minipage}{0.7\textwidth}
\begin{figure}[H]
\begin{centering}
\subfloat[\label{fig:05a}]{\includegraphics[height=4.5cm]{Fig05a}}
\subfloat[\label{fig:05c}]{\includegraphics[height=4.5cm]{Fig05c}}
\end{centering}
\end{figure}
\end{minipage}\hfill \begin{minipage}{0.25\textwidth}
\captionof{figure}{\ref{fig:05a} Abaissement du niveau de Fermi lors d'une oxydation ; \ref{fig:05c} Augmentation du niveau de Fermi lors d'une réduction}
\end{minipage}


En effet, même si nous ne pouvons pas établir d'équivalence, nous admettrons la règle suivante.

\definition{L'énergie de Fermi représente l'énergie potentielle d'un électron dans une phase condensée à \SI{0}{\kelvin} et en l'absence de travail des forces électriques extérieur. Le potentiel chimique de l'électron, introduit au Chp. \ref{chp:thermochimie} est l'analogue thermodynamique de l'énergie de Fermi, et présente exactement le même sens. En réalité, ces grandeurs sont équivalentes à une conversion d'unité près.

\begin{align*}
\epsilon_F ( 0 ~ \si{\kelvin}) = \frac{1}{N_A} \mu_{e^-} (0 ~ \si{\kelvin})
\end{align*}

Dans une cellule électrochimique en circuit ouvert, on pourra alors écrire : 
\begin{equation}
-e_{f.e.m.} = \frac{\epsilon_F^\oplus - \epsilon_F^\ominus}{e}
\end{equation}

où $e$ indique la charge élementaire de l'électron et $e_{f.e.m}$ la force électromotrice.}


\subsection{Autres grandeurs du modèle des bandes}

\subsubsection{Densité d'état}

Dans les pays anglophones, cette grandeur appelée Density-Of-States\mySideNote{Souvent notée DOS dans les exercices, et en général, dans la littérature}, peut se lire de la manière suivante : 
\begin{equation}
\rho_\epsilon ( \epsilon ) \cdot d \epsilon \text{ : Nombre d'état d'énergie }\epsilon\text{ à } d\epsilon \text{ près.}
\end{equation}

Ainsi, la valeur de cette fonction vous indique comme sont réparties les niveaux d'énergie au sein des bandes, et combien d'état sont disponibles à une énergie donnée. Dans la Fig. \ref{fig:05b}, on voit qu'il existe plus d'états électroniques au milieu de la bande qu'à ses limites.

\definition{La densité est une fonction mathématique, notée $\rho_\epsilon (\epsilon)$, qui représente la variation du nombre d'état en fonction de l'énergie. Elle a pour unité l'inverse d'une énergie.}



\subsubsection{Gap d'énergie}

Les bandes d'énergies ne sont pas forcément continues comme nous le verrons par la suite, mais il peut exister une zone où aucun état existe\mySideNote{c'est-à-dire que la densité d'état est nulle} entre la bande de valence et de conduction.

\definition{Lorsque sur l'intervalle $[\epsilon_v, \epsilon_c]$, la densité d'état est nulle sur une plage d'énergie grande devant l'énergie thermique, on parle d'énergie de gap (ou avec l'anglicisme gap d'énergie). L'énergie de gap est définie comme : 
\begin{align*}
E_g = \epsilon_c - \epsilon_v
\end{align*}
et est le plus souvent exprimée en \si{\electronvolt}.}


\section{Comprendre les bandes au sein des oxydes et dérivés}

Les oxydes métalliques sont composés de liaisons iono-covalent, majoritairement ioniques. Lorsqu'un cation et anion forment une liaison ionique comme représenté à la Fig. \ref{fig:02-07c}, on sait que les niveaux de haute (resp. basse) énergie sont centrés sur l'orbitale atomique du cation (resp. anion).

En général, on classifie les bandes en deux catégories dans les oxydes :
\begin{itemize}
\item les \blancbis{bandes anioniques}, qui sont les plus basses en énergie ;
\item les \blancbis{bandes cationiques}, qui sont les plus hautes en énergie.
\end{itemize}

\attention{
Il doit vous sembler naturel que les niveaux des anions soient plus bas en énergie que ceux des cations. Sinon, cela signifierait que les anions s'oxyderaient au profit des cations$\ldots$
}

Le physicien américain J.-B. Goodenough, prix nobel 2019\mySideNote{avec S. Whittingham et A. Yoshiro pour \og leur contribution à la création des batteries au lithium-ion \fg{} . On est pile dans le sujet !}, a établi des schémas de bande qualitatif pour les principaux oxydes simples et complexes. On va se proposer de construire progressivement ces diagrammes en partant de la liaison ionique.

L'ensemble sera illustré sur l'exemple du monoxyde de titane \ce{TiO}, qui cristallise en mode cubique face centré\mySideNote{cf. \ce{CoO} et \ce{NiO} en 3.05}. 

\subsection{Modèle de bande en supposant une liaison purement ionique}

D'après les règles de l'Aufbau, le remplissage électronique donne : 
\begin{align*}
_8[\ce{O}] &= ~ _2[\ce{He}]2s^2 2p^4 \\
_{22}[\ce{Ti}] &= ~ _{18}[\ce{Ar}]4s^2 3d^2 4p^0
\end{align*}

Dans l'oxygène \ce{O} (resp. \ce{Ti}), les orbitales de valences sont $2s$ et $2p$ (resp. $3d$, $4s$ et $4p$) et présentent $6$ (resp. $4$) électrons de valence. On a au final 26 états possibles\mySideNote{donc 13 niveaux d'énergie puisqu'il y a 2 états de spin par niveaux} par motif, dont 10 sont occupés.

Selon la théorie du champ cristallin, l'éclatement de orbitales $d$ est dû à la géométrie et l'intensité de la coordination des métaux de transition aux plus proche voisins (dans un oxyde \ce{O^{2-}}).

Si la liaison est purement ionique\mySideNote{"grand" écart $\Delta \chi$}, on s'attend à avoir des orbitales localisées sur les atomes :

\begin{minipage}{0.48\textwidth}
\centering
\includegraphics[width=0.5\textwidth]{Fig26a}
\end{minipage}\begin{minipage}{0.48\textwidth}
\captionof{figure}{Diagramme d'interaction entre deux orbitales pour une iaison \SI{100}{\percent} ionique entre le métal et l'oxygène.}
\end{minipage}

On obtient le diagramme d'orbitale moléculaire suivant que l'on peut "étendre" au cristal.

\begin{center}
\includegraphics[width=0.95\textwidth]{Fig26b}
\captionof{figure}{Diagramme d'orbitale moléculaire du monoxyde de titane, et extension au diagramme de bande, en supposant une liaison purement ionique}
\end{center}

Les bandes d'énergie dans les liaisons ioniques sont très étroites, on a rarement superposition des bandes. On dit que ce sont des \og{} états localisés \fg{} sur les états.

On retrouve bien que les états de basses (resp. hautes) énergies sont anioniques (resp. cationiques). Le remplissage électronique entre états anioniques et cationiques montre bien que l'on a une liaison ionique : 
\begin{align*}
\ce{Ti-O <-> Ti^{2+} + O^{2-}}
\end{align*}

Ce diagramme ionique type est à avoir en tête, car c'est la description la plus basique que l'on puisse faire des bandes. Les choses se compliquent lorsque l'on ajoute la covalence.

\subsection{Modèle de bande en supposant une liaison iono-covalente, faiblement covalente}

Prendre en compte la covalence revient à dire qu'il existe des recouvrements orbitalaires entre les orbitales, ce qui signifie que lors d'une interaction orbitalaire, les électrons de l'anion (resp. centre métallique) sont partiellement partagés le centre métallique (resp. l'anion).

\begin{minipage}{0.48\textwidth}
\centering
\includegraphics[width=0.5\textwidth]{Fig26c}
\end{minipage}\begin{minipage}{0.48\textwidth}
\captionof{figure}{Diagramme d'interaction entre deux orbitales pour une liaison iono-covalente entre le métal et l'oxygène.}
\end{minipage}

Il existe trois possibilités de recouvrement dans une structure métallique.

\begin{figure}[H]
\centering
\subfloat[\label{fig:27a}]{\includegraphics[height=3cm]{Fig27a}}\hfill\subfloat[\label{fig:27b}]{\includegraphics[height=3cm]{Fig27b}}\hfill\subfloat[\label{fig:27c}]{\includegraphics[height=3cm]{Fig27c}}
\caption{Recouvrement orbitalaires possibles \ref{fig:27a} métal-anion, selon l'axe de la liaison ; \ref{fig:27b} métal-anion, perpendiculairement à l'axe de la liaison et \ref{fig:27c} métal-métal}
\end{figure}

Les recouvrements métal-anion sont les plus fréquemment rencontré et sont soit axial\mySideNote{noté $\sigma$} (voir Fig. \ref{fig:27a}), soit latéral\mySideNote{noté $\pi$} (voir Fig. \ref{fig:27b}). Parfois, il existe des recouvrements métal-métal, qui sont d'autant plus probable : 
\begin{itemize}
\item le centre métallique a un nombre quantique principal grand ;
\item le centre métallique est en début de période ;
\item le degré d'oxydation du métal est faible ;
\item les centres métalliques sont "proches".
\end{itemize}

On admettra que toutes ces interactions donnent un diagramme d'orbitale qui a cette forme : 

\begin{center}
\includegraphics[width=0.95\textwidth]{Fig26d}
\captionof{figure}{Diagramme d'orbitale moléculaire du monoxyde de titane, et extension au diagramme de bande, en supposant une liaison iono-covalente}
\end{center}

La covalence a deux conséquences immédiates : 
\begin{itemize}
\item une diminution $\Delta E_+$ (resp. augmentation $\Delta E_-$) de l'énergie de niveaux liants (resp. antiliants) d'autant plus forte que l'interaction entre les orbitales est grande ;
\item un \blancbis{élargissement des bandes} qui est d'autant plus grand que l'interaction entre les orbitales est grande.
\end{itemize}

\definition{On gardera en tête que l'interaction entre deux orbitales évolue en :
\begin{align*}
\text{interaction} \propto \frac{s^2}{\Delta E}
\end{align*}
où $\Delta E$ est l'écart d'énergie et $s$ le recouvrement entre les deux orbitales qui interagissent.
L'interaction joue sur l'\blancbis{éclatement des bandes} et la \blancbis{largeur de bande}.
}

\section{"Manipuler" les bandes quand on est chimiste}

\subsection{Rédox classique versus rédox interne}

Lorsque l'on charge une batterie, il y a apparition de phénomènes d'oxydation interne entre le(s) métal(ux) de transition et le lithium au pôle $\oplus$. Par exemple, dans un dioxyde métallique lithié \ce{Li_xMO2}: 
\begin{center}
\begin{tabular}{ccc}
$SoC = \SI{100}{\percent}$ & $x=0$ & no(\ce{M}) = +IV \\ 
$SoC = \SI{0}{\percent}$ & $x=1$ & no(\ce{M}) = +III
\end{tabular}
\end{center}

Classiquement, comme représenté à la Fig.\ref{fig:02-08a}, ce sont les orbitales $d$ du métal\mySideNote{orbitales cationiques} qui sont concernées par les processus rédox.
\begin{figure}[H]
\begin{centering}
\subfloat[\label{fig:02-08a}]{\includegraphics[height=5cm]{Fig08a}}\hfill
\subfloat[\label{fig:02-08b}]{\includegraphics[height=5cm]{Fig08b}}\hfill
\subfloat[\label{fig:02-08c}]{\includegraphics[height=5cm]{Fig08c}}
\end{centering}
\caption{Allure du diagramme de bande lorsque \ref{fig:02-08a} les interactions orbitalaires sont faibles ; \ref{fig:02-08b} les interactions orbitalaires sont fortes, causant l'apparition de rédox interne ; \ref{fig:02-08c} le processus de rédox interne a lieu.}
\end{figure}

Cependant, quand les interactions covalentes sont très fortes, comme représenté à la Fig.\ref{fig:02-08b}, les bandes aioniques peuvent être très proche du niveau de Fermi. Dans ce genre de situation, il apparaît des phénomènes de \blancbis{rédox interne} où l'oxydation n'est plus localisée sur le métal, mais aussi sur l'anion.


\subsection{Substitution}

La principale manière qu'a le chimiste de contrôler la position des bandes, c'est de réaliser des substitutions.

Soit la substitution est \blancbis{cationique}, ce qui permettra de jouer sur la position des bandes $d$. 

Soit la substitution est \blancbis{anionique}, ce qui permettra de jouer sur la position des bandes des anions et de manière indirect\mySideNote{interaction orbitalaire} jouer sur la position des bandes cationiques $d$.

Au chapitre suivant, on sera emmené à visiter un cas d'école pour voir comment vous pouvez jouer sur la composition chimique des solides pour obtenir les propriétés désirées.
\end{comment}





\begin{comment}
\definition{

\paragraph*{Règle 1} L'éclatement de orbitales $d$ est dû à la géométrie et l'intensité de la coordination des métaux de transition aux plus proche voisins (\ce{O^{2-}}). Cet éclatement est similaire à celui obtenu par la théorie du champ cristallin.



\paragraph*{Règle 2} Deux orbitales atomiques ne se combinent que si leur recouvrement $S$ est non nul. Pour cela, elles doivent possèder les mêmes propriétés de symétries présenter un faible écart en énergie $\Delta E$.

\paragraph*{Règle 3} L'éclatement entre niveaux liants et anti-liants résultant est d'autant plus grand que l'écart d'énergie $\Delta E$ des orbitales atomiques est petit et que le recouvrement $S$ est grand.

\paragraph*{Règle 4} Du fait de la grande quantité d'interactions inter-atomiques, les niveaux d'énergie s'élargissent en bande dont la position est donnée par :

\begin{itemize}
\item la position des niveaux des atomes isolés ;
\item de la stabilisation par l'environnement chimique (un cation est finalement "solubilisé" dans une mer de \ce{O2^{2-}}).
\end{itemize}

\paragraph*{Règle 5} La largeur des bandes augmente (donc les gaps d'énergie diminuent) lorsque : 
\begin{enumerate}
\item on part des retrouvrements d'orbitales de type $p$ vers des recouvrement de type $s$ ;
\item on passe des orbitales les plus basses vers les orbitales les plus hautes en énergie.
\end{enumerate}


\paragraph*{Règle 6} le nombre total d'état est conservé quand on passe des atomes aux cristals.}
\end{comment}


%\begin{center}
%\includegraphics[height=5cm]{Fig06}
%\captionof{figure}{Fig06}
%\end{center}




\begin{comment}
%\subsection{Aspect électroniques des matériaux d'électrode}


Structure de bande : bande 2p
Potentiel rédox : niveau de Fermi
- niveau occupés = oxydation
- niveau vide = réduction

Aspect 

Chimie quantique du solide : décrire le cristal comme un ensemble de maille successifs

-> Maille Li, Li2, Li3, Li4
-> bande d'énergie (continuum de niveaux électroniques)

Oxydation : Mettre des électrons dans des niveaux anti-liants -> déstabilisation
Réduction : Enlever des électrons -> déstabilisation électronique
-> Interaction entre les atomes de Li faible : 
-> Gap d'énergie : gap Mott-Hubbard (UHP, LHB)

Interaction AB : 
Liaison Ti-O : diagrammde d'énergie intra-maille

Situation : Interaction intra-maille > inter-maille
Situations rédox classiques

Situation : Interaction inter-maille > intra-maille
Transfert d'électron d'une band à une autre

-> Nature des bandes = définie le potentiel de la cellule
-> nombre de bande = contrôle de la capacité

Molécule AB : Liaison ionique vs liaison iono-covalente.
Transfert total d'électron vs transfert partiel d'électron
Rédox atomique ou rédox moléculaire.

Dureté chimique corrélable

%\subsection{Structures cristallographiques}

Classification 0D, 1D, 2D, 3D
Sites Td versus sites Oh

Structure cristallographique : site pour accueillir un ion Li

Formes élémentaires : Td, Oh

Connectivité des polyèdres : 
- sommets, arrête, faces

OD : 
1D : NaCoSO4F.2H2O, MX3, AFeS2, A Mo3X3, nanotubes de C
2D : CoO2, graphite, dichalcogénures, oxyhalogénures
3D : LiFeSO4F, spinelles, phases de chevrel, péroxskites, rutiles

Li+ : sites d'accueil (Td, trigonal prismatique, Oh) , réaction d'inversion : 
Plus la répulsion est élevée : sites octaedriques plutôt que trigonal prismatique.
Ex : TiS2, NbS2 (charge soufre-séleniures diffuse)



V2O5, LiCoO2, Ni(OH)2
-> augmenter le no du métal : V(+V)
-> intercalation du lithium (synthèse in situ avec du lthium)
-> proton

%Structure de LiMn2O4
%https://www.youtube.com/watch?v=2_IILk7jzkQ&ab_channel=MaterialsScienceEducation
%Phase transtion Cuvbic/Tetragonal


Lien : Maille élémentaire -> Cristal



\subsubsection{Influence de la substitution}


Relation entre potentiel d'électrode et structure électronique}


%https://www.youtube.com/watch?v=zdZIwLedZ5w&t=2518s&ab_channel=EliasSebti
%https://www.youtube.com/watch?v=UIwAKPvg4vo&list=PLSIQyXlfw9qceUWcpZPPhNVdxzjIeT7V3&ab_channel=MaterialsScienceEducation

Structure de bande : bande 2p
Potentiel rédox : niveau de Fermi
- niveau occupés = oxydation
- niveau vide = réduction

Aspect 

Chimie quantique du solide : décrire le cristal comme un ensemble de maille successifs

-> Maille Li, Li2, Li3, Li4
-> bande d'énergie (continuum de niveaux électroniques)

Oxydation : Mettre des électrons dans des niveaux anti-liants -> déstabilisation
Réduction : Enlever des électrons -> déstabilisation électronique
-> Interaction entre les atomes de Li faible : 
-> Gap d'énergie : gap Mott-Hubbard (UHP, LHB)

Interaction AB : 
Liaison Ti-O : diagrammde d'énergie intra-maille

Situation : Interaction intra-maille > inter-maille
Situations rédox classiques

Situation : Interaction inter-maille > intra-maille
Transfert d'électron d'une band à une autre

-> Nature des bandes = définie le potentiel de la cellule
-> nombre de bande = contrôle de la capacité

Molécule AB : Liaison ionique vs liaison iono-covalente.
Transfert total d'électron vs transfert partiel d'électron
Rédox atomique ou rédox moléculaire.

Dureté chimique corrélable
\end{comment}